%%%%%%%%%%%%%%%%%%%%%%%%%%%%%%%%%%%%%%%%%%%%%%%%%%%%%%%%%%%%%%%%%%%
\section{Hardware Implications}
\label{sec:Hardware}
%%%%%%%%%%%%%%%%%%%%%%%%%%%%%%%%%%%%%%%%%%%%%%%%%%%%%%%%%%%%%%%%%%%
% (1) the tax as a share of the host's own traffic, datacenter and edge;
% (2) the cache cliff, which is where the edge case is already close;
% (3) where the engine could sit; (4) device candidates; (5) the cost
% model we have not filled in, stated as such.

\subsection{Monitoring as a share of the host's own traffic}
\label{subsec:Tax}
Eq.~(\ref{eq:gputraffic}) is easier to interpret against what the model
itself already moves. In the datacenter, a 70B model at fp8 reads roughly
\qty{70}{\giga\byte} of weights per token; a library of $K=100$K rows at
$d=8{,}192$ adds \qty{1.6}{\giga\byte}, or about 2\%. At the edge, an 8B
model at 4-bit reads roughly \qty{4}{\giga\byte} per token, and $K=25$K
at $d=4{,}096$ adds \qty{205}{\mega\byte}, about 5\%, rising to roughly
20\% by $K=100$K. The same architectural point therefore arrives one to
two orders of magnitude earlier in $K$ at the edge, on parts that have no
HBM available at any price. Edge inference is already bandwidth-limited,
so every monitoring byte trades approximately 1:1 against token rate, and
the 15--\qty{60}{\watt} envelope makes DRAM energy per bit a first-order
term rather than a footnote.

\subsection{The cache cliff, and how close current banks are to it}
\label{subsec:CacheCliff}
Capacity produces a sharper discontinuity than bandwidth. Consider a
Jetson-class part: \qty{273}{\giga\byte\per\second} LPDDR5X, roughly
\qty{4}{\mega\byte} of shared L2 and no HBM~\cite{nvidia_thor}, running
an 8B assistant onboard and monitoring only the four behaviors observed
in the campaign of Sec.~\ref{sec:Introduction}. At $d=4{,}096$, fp16, 500
stored rows is \qty{4.1}{\mega\byte}, which is that shared L2 almost
exactly. Below the cliff the bank is resident and monitoring is close to
free; above it the bank streams from LPDDR on every token.

The arithmetic of Sec.~\ref{sec:Introduction} put four behaviors at
520--1{,}000 rows, and the banks currently in hand are already in that
range: 47 signatures across 13 layers is roughly 600 rows, and the
exemplar bank of Sec.~\ref{sec:Results} is about 4{,}000. The fair
pushback is that at the code widths Sec.~\ref{sec:Results} measured,
today's banks compress back under the cliff -- 4{,}000 rows at 2{,}048
bits is about \qty{1}{\mega\byte}. Two qualifications apply. Compression has been validated for detection, not for the
steering-grade full-precision directions that a correcting monitor would
need. And the cliff does not move as the library grows: at
whole-vocabulary scale even the compressed bank is out of cache
(Fig.~\ref{fig:capacity}). The edge crossing is therefore not purely a
$10^{5}$ hypothetical; the libraries already in hand sit near the
threshold.

\subsection{Where the search engine could sit}
\label{subsec:Placement}
Table~\ref{tab:placement} lists the placements and what each trades. The
trade is always the same: how far the query has to travel, and what it
has to contend with on the way. A monitor co-resident on the inference
GPU competes with live inference for cache and bandwidth, which is the
contention the analysis above assumes away and which would need to be
measured. A back-end-of-line (BEOL) search engine in the metal layers
above GPU compute keeps the query travel to microns and the library
nonvolatile and static, so it augments a part already purchased -- but no
foundry PDK offers BEOL \fefet today, and BEOL RRAM/MRAM at 22/16\,nm is
the defensible near-term citation
(the BEOL device work is in Sec.~\ref{subsec:CAMspace}), which makes
that row a 2030s option and it should be labeled that way.
\dt[confirm you want the BEOL row in a DATE paper at all. It is the most
speculative entry in the table and DATE reviewers are less forgiving
than a sponsor audience. The CXL/appliance row is the one an industrial
reviewer can price today]{}

\begin{table}[t]
\centering
\caption{Placements for a \cam-resident monitor.}
\label{tab:placement}
{\footnotesize
\begin{tabular}{@{}p{0.27\columnwidth}p{0.66\columnwidth}@{}}
\toprule
\textbf{Placement} & \textbf{Trade-off} \\
\midrule
Co-resident on the inference GPU &
Contends with live inference for cache and bandwidth; the contention
term is not in our numbers and could dominate at scale. \\
Dedicated monitoring GPU(s) &
Whole GPUs scanning a static table, plus cross-device latency. \\
\cam or CXL-attached appliance &
Matches in place, but the query crosses a board-level bus; a separate
part to build and qualify. \\
BEOL search engine above GPU compute &
Query travels microns, no cache contention, nonvolatile static library;
no PDK support for BEOL \fefet today. \\
\bottomrule
\end{tabular}}
\end{table}

\subsection{Device candidates}
\label{subsec:Devices}
Table~\ref{tab:devices} lists candidate substrates. The encodings of
Sec.~\ref{sec:Mapping} matter here because they change which row is
required. A thermometer code needs only the last row, i.e., no new
device: a binary SRAM \cam in any CMOS process. The \signrp ternary code
needs wildcard capability, available in CMOS \tcam at higher area or in a
2-transistor \fefet cell~\cite{ni2019fetcam}. A multi-bit Euclidean
matchline is the row that requires a multi-level
device~\cite{kazemi2021fefetcam}. The two-stage variant would benefit
from it, as would the first pipeline of Sec.~\ref{subsec:FirstPipeline},
whose 3-bit Euclidean matching maps directly onto such a cell, and the
frozen-projection Euclidean design of Sec.~\ref{subsec:PathBack}. The
multi-bit row therefore stays in the table even though the encodings
this paper measures do not require it: it is the substrate for the next
experiment rather than a device the current results depend on.

\textbf{The two workloads may want different cells:} The measured
results give the multi-bit row a second reason to stay. At a matched
384 bits the \signrp ternary code is ahead of thermometer on the
deception grid (Sec.~\ref{sec:Results}), while on the readout workload
the 3-bit Euclidean code is well ahead of every Hamming-metric encoding
at the same storage budget (Sec.~\ref{subsec:JspaceFidelity}). The two
workloads are asking different questions -- separating two labeled
classes there, reproducing a specific ranked list here -- and the
ordering flip is consistent with the readout needing the magnitude
information that sign codes discard. Two of the checks on this are
now in: at matched cell count the deception-grid ordering is within
noise (Sec.~\ref{subsec:Grids}), and a common frozen projection on
the deception grid removes the Euclidean advantage there
(Sec.~\ref{subsec:BitWidth}), so the flip is a property of the
readout workload and its trained projection rather than of the cell.
Whether it survives a common frozen projection on the readout workload
itself is the remaining check (Sec.~\ref{subsec:NextMeasurements}).
If it does, a monitoring engine that has to serve both ends of
Fig.~\ref{fig:ladder} either provisions two cell types or pays for the
multi-bit cell and serves the small-$K$ workload with it too; if not,
one array serves both.\ib{Items 1 and 8(a) feed this paragraph; the
readout-side common-projection check is still open.} Write endurance appears in
the table because the bank is rewritten whenever the monitored model is
updated or the probe layer is re-tuned, which is a per-model-release
event rather than a per-token one; that is the property that makes an NVM
bank reasonable here and unreasonable for transformer intermediates.

\begin{table}[t]
\centering
\caption{Candidate substrates for the signature bank. Device figures are
literature values, not our measurements.}
\label{tab:devices}
{\footnotesize
\begin{tabular}{@{}lcccl@{}}
\toprule
\textbf{Tech} & \textbf{Bits/cell} & \textbf{Write} & \textbf{Endurance} &
\textbf{Best fit here} \\
\midrule
\fefet   & 2--3   & $\sim$20\,ns    & $>$5$\times$10$^{10}$ & multi-bit Euclidean \\
ECRAM    & up to 6 & $\sim$200\,$\mu$s & $\sim$10$^{12}$    & static direction bank \\
RRAM     & 1--3   & sub-$\mu$s      & $\sim$10$^{9}$        & density; wear-leveling \\
SRAM \cam & 1     & SRAM            & --                    & Hamming; any CMOS fab \\
\bottomrule
\end{tabular}}
\end{table}

\subsection{A first-order energy and latency estimate}
\label{subsec:CostModelGap}
The analysis above is capacity and bandwidth arithmetic. What can be
stated from published device and memory figures is a first-order
estimate of the energy and latency of the search itself; a measured
GPU baseline is reported below it, and the gap between the two sides
is wide enough to survive considerable error in the inputs.

On the GPU side, the dominant energy term at large $K$ is reading the
bank out of DRAM. Measured HBM2 access energy is approximately
3.97~pJ/bit~\cite{oconnor2017fgdram}. At the whole-vocabulary
operating point -- $K=1.3\!\times\!10^{5}$ rows of fp16 at $d=8{,}192$,
i.e., the \qty{2.1}{\giga\byte} bank of
Sec.~\ref{subsec:ReadoutCost} -- reading the bank once is roughly
\emph{67~mJ per token per layer}, before any arithmetic is charged. On
the \cam side, the same search is $K$ rows of 2{,}048-bit code compared
in place, about $2.6\!\times\!10^{8}$ cell comparisons. Published
figures for binary and ternary \cam search sit in the range of
0.1--1~fJ per cell per search across CMOS survey data and fabricated
\fefet arrays~\cite{pagiamtzis2006cam,ni2019fetcam,liu2022evacam},
which puts the whole-bank search at roughly \emph{0.03--0.3~$\mu$J per
token per layer} -- five to six orders of magnitude under the DRAM
read. The gap decomposes into two factors: $64\times$ from the code
width (2{,}048 bits stored per row against 131{,}072; measured
fidelity at that compression is what Secs.~\ref{sec:Results}
and~\ref{sec:LargeK} establish) and $\sim$$4\!\times\!10^{3}$ from the
substrate (fJ-scale cell comparison against pJ-scale DRAM bit
movement). As such, even a $10\times$ error in the per-cell figure --
which spans the published range -- moves the conclusion by one order
out of five or six. Scaled to a deployment, the same arithmetic at 25
monitored layers and $10^{3}$~tokens/s is roughly 1.7~kW of DRAM
access power on the GPU side against tens of milliwatts of array
search: monitoring at this $K$ is not a workload a deployed part
absorbs, which is what Fig.~\ref{fig:ladder}(b) says in bandwidth
terms.\footnote{What the estimate omits, on each side: the GPU figure
excludes the distance arithmetic itself and any cache reuse across
batched requests (Sec.~\ref{subsec:NextMeasurements}); the \cam figure
takes array-level published values, which include matchline and
sense-amp energy but not the query projection ($\sim$$10^{6}$
multiply--accumulates per token, negligible at this scale) or the
mismatch-count reduction across subarrays. These are first-order
literature values, not our measurements.}

Latency divides the same way. Streaming \qty{2.1}{\giga\byte} at an
H100's \qty{3.35}{\tera\byte\per\second} is a 0.63~ms floor per token
per layer at batch~1, or roughly 10~$\mu$s per request amortized at
batch 64. Fabricated and modeled TCAM arrays report search times of
0.3--4.4~ns~\cite{liu2022evacam,ni2019fetcam}; with subarrays searched
in parallel and mismatch counts reduced in a tree, a whole-bank search
plausibly completes in tens of nanoseconds, and even at 100~ns it sits
two orders under the amortized floor and roughly four under the
batch-1 floor.

\textbf{A measured GPU baseline:} On an H100 PCIe, scanning the raw
fp16 bank at 131K rows with the $k=31$ top-$k$ included and board
power sampled during the timed loop costs 1.22~ms and 277~mJ per token
at batch~1, and 21~$\mu$s and 4.7~mJ per token at batch~64. Achieved
bandwidth was 1.6--1.8~TB/s, about half of the part's specification,
so the $K\!\approx\!204$K crossing of Sec.~\ref{sec:Scaling} is
closer to 107K in practice. Compression also helps the GPU, and the
comparison should say so: the same 384-bit codes the \cam holds,
searched on the GPU with a population-count kernel (FAISS), cost
28~$\mu$s and 2.1~mJ per token at batch~1 (0.45~mJ above idle) and
0.5~$\mu$s and 0.04~mJ at batch~64, flat in $K$ from $10^{3}$ to 131K
rows -- the scan itself is under 30~$\mu$s, and the GPU sits at about
3\% of its per-token budget at $10^{3}$~tokens/s. What the array
offers over that baseline is the search energy itself
(0.03--0.3~$\mu$J against 40~$\mu$J--2~mJ per token per layer) and
freedom from contention with live inference, rather than bandwidth
relief alone. Table~\ref{tab:costmodel} carries these entries, the
\cam rows marked as estimates, alongside the rows that remain
open.\ib{Item 8(b): two H100 PCIe baselines, each including the $k=31$
top-$k$, at $K$ up to 131K and batch 1 and 64, board power sampled
during the timed loop. Framing of the popcount result is for MN to
confirm.}

\begin{table}[t]
\centering
\caption{The cost model: measured rows, first-order estimates, and open
entries. \textsc{est} marks first-order values from literature figures
(HBM at 3.97~pJ/bit~\cite{oconnor2017fgdram}; \cam cells at
0.1--1~fJ/search~\cite{pagiamtzis2006cam,ni2019fetcam,liu2022evacam}),
at $K=1.3\!\times\!10^{5}$, 2{,}048-bit rows, per token per layer.
\textsc{meas} marks measurements on one H100 PCIe at batch 1 / batch
64; \textsc{tbd} entries require the measurements of
Sec.~\ref{subsec:NextMeasurements}.}
\label{tab:costmodel}
{\scriptsize\setlength{\tabcolsep}{3pt}
\begin{tabular}{@{}lccp{0.23\columnwidth}@{}}
\toprule
\textbf{Metric} & \textbf{CAM} & \textbf{GPU} & \textbf{Assumptions} \\
\midrule
Energy / search & \makecell[l]{0.03--0.3\,$\mu$J\\\textsc{est}} &
  \makecell[l]{277 / 4.7\,mJ \textsc{meas}\\($\sim$67\,mJ \textsc{est})} &
  cell figure spans published range; \textsc{est} is DRAM access only \\
Latency / search & \makecell[l]{0.3--4.4\,ns/arr.\\\textsc{est}} &
  \makecell[l]{1.22\,ms / 21\,$\mu$s \textsc{meas}\\(0.63\,ms \textsc{est})} &
  reduction net.\ open; \textsc{est} is the batch-1 streaming floor \\
GPU on 384-b codes & n/a & \makecell[l]{2.1 / 0.04\,mJ; 28 / 0.5\,$\mu$s\\\textsc{meas}} &
  FAISS popcount, flat in $K$ from $10^{3}$ to 131K \\
Area / $10^{3}$ rows & \textsc{tbd} & n/a &
  cell type (Table~\ref{tab:camspace}), node \\
Coverage / mm$^{2}$, / W & \textsc{tbd} & \textsc{tbd} &
  behaviors $\times$ layers, at fixed detection rate \\
Traffic / token & $d b L_{\mathrm{mon}}$ & $K d b L_{\mathrm{mon}}$ &
  measured here; Eqs.~(\ref{eq:gputraffic}), (\ref{eq:camtraffic}) \\
Bank capacity & \multicolumn{2}{c}{Fig.~\ref{fig:capacity}} &
  measured here; code widths from Sec.~\ref{sec:Results} \\
\bottomrule
\end{tabular}}
\end{table}

\subsection{The measurements that would replace the estimates}
\label{subsec:NextMeasurements}
Each estimate above maps to a specific, short measurement, and the
tooling for the array side already exists in-house.
{\bf (1)} \textbf{Array-level characterization with Eva-CAM:} Eva-CAM
is a circuit/architecture-level evaluation tool for TCAM, MCAM and
ACAM arrays across NVM technologies, validated against fabricated
chips to within 6--20\%~\cite{liu2022evacam}. Running it at a stated
node for exactly the configurations this paper measures -- 384- and
2{,}048-bit rows, banks of $10^{3}$--$1.3\!\times\!10^{5}$ rows, BCAM,
\tcam and 3-bit \mcam cells -- would replace the \textsc{est} entries
of Table~\ref{tab:costmodel} with modeled energy, latency and area,
including the subarray split and the reduction network the footnote
above leaves open. This is days of work with an existing tool.
{\bf (2)} \textbf{A tuned GPU baseline:} the H100 measurements of
Sec.~\ref{subsec:CostModelGap} (fp16 scan and FAISS
popcount~\cite{faisscuvs2025}, batch 1 and 64, measured power) are now
in Table~\ref{tab:costmodel}; the Jetson-class part remains to be
measured, so the amortization a datacenter part enjoys -- and an edge
part does not -- can appear in the table rather than in a caveat.\ib{Item
8(b) done for the H100; Jetson not measured.} Recall
for any approximate index should additionally be measured under
queries chosen \emph{against} the index rather than drawn at random,
since that is the regime a monitor operates
in~\cite{andoni2026robustann}.
{\bf (3)} \textbf{The fidelity check the encoding comparison still
needs:} matched cell count (Sec.~\ref{subsec:Grids}) and a common
frozen projection on the deception grid (Sec.~\ref{subsec:BitWidth})
are done; a common projection across the Hamming and Euclidean
encodings on the readout workload remains, and it is what settles
whether the cell-type flip of Sec.~\ref{subsec:Devices} is a device
result or a projection artifact. It is a cluster run on data already
in hand.\ib{Items 1 and 8(a) done; the readout-side run is not.}
\dt[items (1) and (2) are the preemptive strike for a reviewer asking
``is the CAM actually cheaper'': the paper now has the first-order
numbers, and this list is the credible path to measured ones. If there
is student time before the deadline, (1) is the highest
value-per-day]{}
